\begin{trapbox}

	\begin{description}[style=nextline, leftmargin=2.8em, labelsep=0.5em, itemsep=0.35em]
		\item[\textbf{\textcolor{CTrap}{易错点一（焦点对应）}}]
			计算距离比时，将右焦点与左准线搭配，或混用焦点准线对。
		\item[\textbf{\textcolor{CTrap}{正确理解（对应规则）：}}]
			每个焦点只对应同侧的准线：右焦点对应右准线 $x=a^{2}/c$，左焦点对应左准线 $x=-a^{2}/c$。
		\item[\textbf{\textcolor{CTrap}{错因分析（记忆偏差）：}}]
			未注意圆锥曲线中焦点与准线的位置对应关系，记忆模糊。
	\end{description}
	\par

	\medskip
	\begin{description}[style=nextline, leftmargin=2.8em, labelsep=0.5em, itemsep=0.35em]
		\item[\textbf{\textcolor{CTrap}{易错点二（绝对值处理）}}]
			在计算点到准线距离时，直接使用 $x-\frac{a^{2}}{c}$ 而不取绝对值，导致符号错误。
		\item[\textbf{\textcolor{CTrap}{正确理解（距离非负）：}}]
			点到直线的距离公式应使用绝对值，即 $|PH|=|x-\frac{a^{2}}{c}|$（或 $|y-\frac{a^{2}}{c}|$ 当焦点在 $y$ 轴时）。
		\item[\textbf{\textcolor{CTrap}{错因分析（代数疏忽）：}}]
			将代数式直接当作距离，忽略了距离必须非负，且点的横坐标可能小于准线坐标。
	\end{description}
	\par

	\medskip
	\begin{description}[style=nextline, leftmargin=2.8em, labelsep=0.5em, itemsep=0.35em]
		\item[\textbf{\textcolor{CTrap}{易错点三（离心率范围）}}]
			误认为只要距离比是常数 $k$ 就一定是双曲线，忽略了 $k$ 是否大于 $1$。
		\item[\textbf{\textcolor{CTrap}{正确理解（范围判定）：}}]
			圆锥曲线统一定义中，$k>1$ 时为双曲线；$0<k<1$ 时为椭圆；$k=1$ 时为抛物线。必须根据 $k$ 判断曲线类型。
		\item[\textbf{\textcolor{CTrap}{错因分析（概念混淆）：}}]
			只记住了“比值等于离心率”，但忘记了椭圆、双曲线、抛物线的离心率范围不同。
	\end{description}

\end{trapbox}
